% Bagian: incompleteness
% Bab: representability-in-q
% Subbagian: c-representable

\documentclass[../../../include/open-logic-section]{subfiles}

\begin{document}

\olfileid[id]{inc}{req}{cre}
\olsection{Fungsi-Fungsi dalam $C$ Dapat Direpresentasikan dalam $\Th{Q}$}

Kita harus menunjukkan bahwa setiap fungsi dalam $C$ dapat direpresentasikan
dalam $\Th{Q}$. Pada akhirnya, kita perlu menunjukkan cara menetapkan kepada
setiap fungsi berargumen $k$, $f(x_0,\dots,x_{k-1})$, dalam~$C$ suatu formula
$!A_f(x_0,\dots,x_{k-1},y)$ yang merepresentasikannya.

\begin{lem}
Diberikan bilangan asli $n$ dan $m$, jika $n \neq m$, maka $Q \Proves \num
n \neq \num m$.
\end{lem}

\begin{proof}
Gunakan induksi pada $n$ untuk menunjukkan bahwa, bagi setiap $m$, jika
$n \neq m$, maka $Q \Proves \eq/[\num n][\num m]$.

Dalam kasus dasar, $n = 0$. Jika $m$ tidak sama dengan $0$, maka
$m = k + 1$ untuk suatu bilangan asli~$k$. Kita mempunyai aksioma yang
menyatakan $\lforall[x][\eq/[0][x']]$. Menurut suatu aksioma kuantor, dengan
mengganti $x$ dengan $\num k$, kita dapat menyimpulkan
$\eq/[0][\num k']$. Namun, $\num k'$ tidak lain ialah $\num m$.

Dalam langkah induksi, kita dapat mengasumsikan bahwa klaim tersebut benar
untuk $n$, lalu mempertimbangkan $n+1$. Misalkan $m$ sebarang bilangan asli.
Ada dua kemungkinan: $m = 0$, atau terdapat suatu $k$ sedemikian sehingga
$m = k+1$. Kasus pertama ditangani seperti di atas. Dalam kasus kedua,
andaikan $n+1 \neq k+1$. Maka $n \neq k$. Menurut hipotesis induksi untuk
$n$, kita mempunyai $\Th{Q} \Proves \eq/[\num n][\num k]$. Kita mempunyai
aksioma yang menyatakan
$\lforall[x][\lforall[y][\eq[x'][y'] \lif \eq[x][y]]]$. Dengan menggunakan
aksioma kuantor, kita mempunyai
$\eq[\num n'][\num k'] \lif \eq[\num n][\num k]$. Dengan menggunakan
logika proposisional, dalam $\Th{Q}$ kita dapat menyimpulkan
$\eq/[\num n][\num k] \lif \eq/[\num n'][\num k']$. Dengan modus ponens,
kita dapat menyimpulkan $\eq/[\num n'][\num k']$, sebagaimana diinginkan,
karena $\num k'$ ialah $\num m$.
\end{proof}

Perhatikan bahwa lema tersebut tidak mengatakan banyak hal: pada intinya,
lema itu mengatakan bahwa $\Th{Q}$ dapat membuktikan bahwa numeral-numeral
yang berbeda menunjuk objek yang berbeda. Misalnya, $\Th{Q}$ membuktikan
$0'' \neq 0'''$. Namun, menunjukkan bahwa hal ini berlaku secara umum
memerlukan kehati-hatian. Perhatikan pula bahwa meskipun kita menggunakan
induksi, induksi itu dilakukan \emph{di luar} $\Th{Q}$.

Kita akan dapat merepresentasikan nol, penerus, tambah, kali, fungsi
karakteristik kesamaan, dan proyeksi. Dalam setiap kasus, formula yang sesuai
untuk merepresentasikan fungsi tersebut sepenuhnya langsung; misalnya, nol
direpresentasikan oleh formula
\[
y = 0,
\]
penerus direpresentasikan oleh formula
\[
x_0' = y,
\]
dan penjumlahan direpresentasikan oleh formula
\[
x_0 + x_1 = y.
\]
Pekerjaannya terletak pada menunjukkan bahwa $\Th{Q}$ dapat membuktikan
pernyataan-pernyataan yang relevan; misalnya, mengatakan bahwa penjumlahan
direpresentasikan oleh formula di atas berarti menunjukkan bahwa untuk setiap
pasangan bilangan asli $m$ dan $n$, $\Th{Q}$ membuktikan
\[
\eq[\num n + \num m][\num {n+m}]
\]
dan
\[
\lforall[y][(\eq[(\num n + \num m)][y] \lif \eq[y][\num {n+m}])].
\]

Bagaimana dengan komposisi? Andaikan $h$ didefinisikan oleh
\[
h(x_0,\dots,x_{l-1}) = f(g_0(x_0,\dots,x_{l-1}), \dots,
g_{k-1}(x_0,\dots,x_{l-1})).
\]
dengan kita telah menemukan formula $!A_f, !A_{g_0}, \dots,
!A_{g_{k-1}}$ yang masing-masing merepresentasikan fungsi
$f,g_0,\dots,g_{k-1}$. Maka kita dapat mendefinisikan formula~$!A_h$ yang
merepresentasikan~$h$ dengan mendefinisikan
$!A_h(x_0,\dots,x_{l-1},y)$ sebagai
\begin{multline*}
  \lexists[z_0\dots][\lexists[z_{k-1}][(!A_{g_0}(x_0,\dots,x_{l-1},z_0) \land
  \dots \land !A_{g_{k-1}}(x_0,\dots,x_{l-1},z_{k-1}) \land]]\\
  !A_f(z_0,\dots,z_{k-1},y)).
\end{multline*}

Terakhir, mari kita perhatikan pencarian tak berbatas. Andaikan
$g(x,\vec z)$ bersifat reguler dan direpresentasikan dalam $\Th{Q}$,
misalnya oleh formula $!A_g(x,\vec z,y)$. Misalkan $f$ didefinisikan oleh
$f(\vec z) = \mu x \; g(x,\vec z)$. Kita ingin menemukan formula
$!A_f(\vec z,y)$ yang merepresentasikan~$f$. Berikut merupakan pilihan yang
alami:
\[
!A_f(\vec z,y) \ident !A_g(y,\vec z,0) \land \lforall[w][(w < y
\lif \lnot !A_g(w,\vec z,0))].
\]

Hal ini dapat ditunjukkan berhasil dengan menggunakan beberapa lema tambahan,
misalnya:

\begin{lem}
  Bagi setiap variabel~$x$ dan setiap bilangan asli~$n$, $\Th{Q}$ membuktikan
  $\eq[(x' + \num n)][(x + \num n)']$.
\end{lem}

Perlu kembali disebutkan bahwa hal ini lebih lemah daripada mengatakan bahwa
$\Th{Q}$ membuktikan
$\lforall[x][\lforall[y][(x' + y) = (x +y)']]$ (yang salah).

\begin{proof}
Seperti biasa, pembuktian dilakukan dengan induksi pada $n$. Dalam kasus
dasar, $n = 0$, kita perlu menunjukkan bahwa $\Th{Q}$ membuktikan
$\eq[(x'+0)][(x+0)']$. Namun, kita mempunyai:
\begin{eqnarray*}
& & \eq[(x' + 0)][x'] \quad \text{menurut aksioma 4} \\
& & \eq[(x + 0)][x] \quad \text{menurut aksioma 4} \\
& & \eq[(x+0)'][x'] \quad \text{menurut baris 2} \\
& & \eq[(x' + 0)][(x + 0)'] \quad \text{menurut baris 1 dan 3}
\end{eqnarray*}
Dalam langkah induksi, kita dapat mengasumsikan bahwa kita telah menurunkan
$\eq[(x' + \num n)][(x + \num n)']$ dalam $\Th{Q}$. Karena $\num{n+1}$
ialah $\num n'$, kita perlu menunjukkan bahwa $\Th{Q}$ membuktikan
$\eq[(x' + \num n')][(x + \num n')']$. Rantai kesamaan berikut dapat
diturunkan dalam $\Th{Q}$:
\begin{eqnarray*}
(x' + \num n') & \eq & (x' + \num n)' \quad \text{aksioma 5} \\
& \eq & (x + \num n')' \quad \text{menurut hipotesis induksi}
\end{eqnarray*}
\end{proof}

\begin{lem}
\begin{enumerate}
\item $\Th{Q}$ membuktikan $\lnot (x < \num 0)$.
\item Bagi setiap bilangan asli~$n$, $\Th{Q}$ membuktikan
\[
x < \num {n+1} \lif (\eq[x][\num 0] \lor \dots \lor \eq[x][\num n]).
\]
\end{enumerate}
\end{lem}

\begin{proof}
Mari kita kerjakan butir~1 dan sebagian butir~2 secara informal (yakni hanya
memberi petunjuk tentang cara menyusun !!{derivation} formalnya).

Untuk butir~1, menurut definisi $<$, kita perlu membuktikan
$\lnot \lexists[y][\eq[(y'+x)][0]]$ dalam $\Th{Q}$, yang ekuivalen
(dengan menggunakan aksioma dan kaidah logika orde pertama) dengan
$\lforall[y][\eq/[(y' + x)][0]]$. Berikut gagasannya: andaikan
$\eq[(y'+x)][0]$. Jika $x$ ialah~0, kita mempunyai
$\eq[(y' + 0)][0]$. Namun, menurut aksioma~4 dari $\Th{Q}$, kita mempunyai
$\eq[(y' + 0)][y']$, dan menurut aksioma~2 kita mempunyai $\eq/[y'][0]$,
suatu kontradiksi. Jadi, $\lforall[y][\eq/[(y' +x)][0]]$. Jika $x$ bukan
$0$, menurut aksioma~3 terdapat suatu~$z$ sedemikian sehingga
$\eq[x][z']$. Namun, kemudian kita mempunyai $\eq[(y' + z')][0]$.
Menurut aksioma~5, kita mempunyai $\eq[(y' + z)'][0]$, yang kembali
bertentangan dengan aksioma~2.

Untuk butir~2, gunakan induksi pada $n$. Mari kita perhatikan kasus dasar
ketika $n =0$. Dalam kasus itu, kita perlu menunjukkan
$x < \num 1 \lif \eq[x][\num 0]$. Andaikan $x < \num 1$. Maka, menurut
aksioma yang mendefinisikan $<$, kita mempunyai
$\lexists[y][\eq[(y'+x)][0']]$. Andaikan $y$ mempunyai sifat tersebut;
jadi, kita mempunyai $\eq[y'+x][0']$.

Kita perlu menunjukkan $\eq[x][0]$. Menurut aksioma~3, jika $x$ bukan~0,
variabel itu sama dengan $z'$ untuk suatu~$z$. Maka kita mempunyai
$\eq[(y' + z')][0']$. Menurut aksioma~5 dari $\Th{Q}$, kita mempunyai
$\eq[(y'+z)'][0']$. Menurut aksioma~1, kita mempunyai
$\eq[(y' + z)][0]$. Namun, menurut definisi, ini berarti $z < 0$,
bertentangan dengan butir~1.
\end{proof}

\begin{explain}
Kita telah menunjukkan bahwa himpunan fungsi yang dapat dihitung dapat
dikarakterisasi sebagai himpunan fungsi yang dapat direpresentasikan dalam
$\Th{Q}$. Bahkan, pembuktian tersebut lebih umum. Dari definisi
kerepresentasian, tidak sulit melihat bahwa setiap teori yang memperluas
$\Th{Q}$ (atau yang memungkinkan $\Th{Q}$ ditafsirkan di dalamnya) dapat
merepresentasikan fungsi-fungsi yang dapat dihitung; sebaliknya, dalam setiap
sistem !!{derivation} yang pengertian buktinya dapat dihitung, setiap fungsi
yang dapat direpresentasikan juga dapat dihitung. Jadi, sebagai contoh,
himpunan fungsi yang dapat dihitung dapat dicirikan sebagai himpunan fungsi
yang direpresentasikan dalam aritmetika Peano, atau bahkan teori himpunan
Zermelo--Fraenkel. Seperti dicatat G\"odel, hal ini agak mengejutkan. Kita
akan melihat bahwa, dalam hal keterbuktian, pertanyaan-pertanyaan sangat peka
terhadap teori mana yang dipertimbangkan; secara kasar, semakin kuat
aksiomanya, semakin banyak yang dapat dibuktikan. Namun, pada cakupan luas
teori-teori aksiomatik, fungsi yang dapat direpresentasikan tepat merupakan
fungsi yang dapat dihitung.
\end{explain}

\end{document}
